\section{Analiza porównawcza rozwiązań europejskich}
\label{sec:porownanie}

\subsection{Przegląd wybranych systemów europejskich}
\label{sec:porownanie-europa}

\subsubsection{Sistema di Interscambio (SdI)}
\label{sec:wstep-sdi}

\textbf{Sistema di Interscambio (SdI)} to krajowy system e-fakturowania obowiązujący we Włoszech, prowadzony przez włoską administrację skarbową. Podobnie do KSeF, oparty jest na modelu \textit{central hub} -- działa jako pośrednik, który weryfikuje, przetwarza i dostarcza faktury. SdI jest najdłużej działającym obowiązkowym systemem B2G w Unii Europejskiej (od 2014~r.) i pierwszym, który objął powszechnym obowiązkiem transakcje B2B i B2C (od 1~stycznia 2019~r.). \cite{edicom-einvoicing-italy,sdi-fatturapa-cose,ec-einvoicing-italy-2025}

\subsubsection{PEPPOL}
\label{sec:wstep-peppol}

\textbf{PEPPOL} (\textit{Pan-European Public Procurement On-Line}) to otwarta, wielostronna sieć interoperacyjna zarządzana przez organizację non-profit \textbf{OpenPeppol AISBL} z siedzibą w Brukseli. W odróżnieniu od SdI i KSeF, PEPPOL nie jest systemem centralnym -- nie przechowuje faktur ani nie pełni roli scentralizowanego pośrednika w obrocie; jest natomiast zbiorem standardów, protokołów i polityk umożliwiających bezpośrednią wymianę ustrukturyzowanych dokumentów między akredytowanymi punktami dostępu (\textit{Access Points}). Sieć obejmuje ponad 40~jurysdykcji na świecie, a jej pierwotny mandat prawny -- obowiązek e-fakturowania B2G -- wynika z~Dyrektywy 2014/55/UE. \cite{peppol-interop-framework,dir-2014-55-ue}

\subsection{Modele architektury: centralizacja vs sieci rozproszone}
\label{sec:porownanie-modele-architektury}

Architekturę wybranych systemów e-fakturowania można podzielić na dwa podstawowe paradygmaty: \textit{model scentralizowany} oraz \textit{model sieci interoperacyjnej}. \cite{edicom-einvoice-models}

\subsubsection{Centralny hub -- model SdI i KSeF}
\label{sec:porownanie-central-hub}

W modelu scentralizowanym każda faktura musi zostać zatwierdzona przez system państwowy, zanim dotrze do odbiorcy. Model ten jest realizowany zarówno przez KSeF oraz SdI, z różniąc się w aspektach takich jak archiwizacja faktur, zakres obowiązku fakturowania oraz potwierdzenie dostarczenia faktury.

\subsubsection{Sieć interoperacyjna -- model 4-corner PEPPOL}
\label{sec:porownanie-4corner}

Sieć PEPPOL opiera się na \textbf{modelu czterech narożników (4-corner model)}, w którym wystawca i odbiorca komunikują się pośrednio przez punkty dostępu (\textit{Access Points}), czyli akredytowanych przez PEPPOL dostawców usług.

\begin{enumerate}
    \item \textbf{C1 -- Nadawca:} wystawca faktury.
    \item \textbf{C2 -- Punkt dostępu nadawcy:} konwertuje fakturę do formatu PEPPOL, waliduje, szyfruje i przesyła do AP odbiorcy, korzystając z dynamicznego wyszukiwania.
    \item \textbf{C3 -- Punkt dostępu odbiorcy:} weryfikuje, deszyfruje, waliduje i przekazuje fakturę do C4.
    \item \textbf{C4 -- Odbiorca:} adresat faktury.
\end{enumerate}

Kluczową różnicą wobec KSeF i SdI jest brak centralnego repozytorium: żaden organ państwowy nie ma automatycznego dostępu do treści faktury w momencie jej transmisji, chyba że sam uczestniczy w wymianie. Wyjątkiem jest \textbf{wariant 5-corner}, będący rozszerzeniem sieci, w którym równolegle do procesu dołączony jest piąty uczestnik (zazwyczaj organ podatkowy) otrzymujący dane z faktury do celów sprawozdawczych. \cite{peppol-interop-framework,einvoicing-dicovering-pepol}

% todo: redakcja

\subsection{Standardy i formaty danych}
\label{sec:porownanie-standardy-formaty}

\subsubsection{Norma europejska EN~16931}
\label{sec:porownanie-en16931}

\textbf{EN~16931} (formalnie EN~16931-1:2017, opublikowana przez Komitet Europejski Normalizacji) to europejska norma treści e-faktury, wymagana Dyrektywą~2014/55/UE dla wszystkich systemów B2G w~Unii. Norma definiuje ok.~180 semantycznych elementów danych (\textit{Business Terms} BT-1\ldots BT-180+) i jest \textit{syntax-agnostic} -- nie narzuca konkretnego formatu pliku. Implementacje zgodne z~EN~16931 mogą używać różnych składni XML, z których dwie zostały oficjalnie zatwierdzone przez Komisję Europejską: \textbf{UBL 2.1} oraz \textbf{UN/CEFACT CII D16B}. \cite{en16931-1,dir-2014-55-ue}

\subsubsection{FatturaPA}
\label{sec:porownanie-fatturapa}

\textbf{FatturaPA} to krajowy schemat XML stosowany w SdI. Dokument XML zbudowany jest wokół elementu głównego \texttt{FatturaElettronica}, zawierającego sekcje \texttt{FatturaElettronicaHeader} (dane nadawcy, odbiorcy, numery referencyjne) i \texttt{FatturaElettronicaBody} (treść faktury, pozycje, kwoty VAT). Atrybut \texttt{FormatoTrasmissione} wskazuje typ faktury: \texttt{FPA12} dla administracji publicznej lub \texttt{FPR12} dla podmiotów prywatnych. Od 1~kwietnia 2025~r. obowiązuje wersja~1.9 specyfikacji. \cite{sdi-fatturapa-format,ec-einvoicing-italy-2025}

FatturaPA nie jest natywnie zgodna z~EN~16931 -- spełnia wymagania merytoryczne Dyrektywy jedynie w~sposób równoważny, ale nie przez implementację normy. Jest to analogia do polskiego FA(3): oba formaty to krajowe schematy XML definiowane autonomicznie przez administracje skarbowe, niezależnie od europejskiej normalizacji.

\subsubsection{PEPPOL BIS Billing 3.0}
\label{sec:porownanie-bis-billing}

\textbf{PEPPOL BIS Billing 3.0} (\textit{Business Interoperability Specification}) to specyfikacja fakturowania stosowana w~sieci PEPPOL, zbudowana jako \textbf{CIUS} (\textit{Core Invoice Usage Specification}) bezpośrednio na normie EN~16931. Oznacza to, że każda faktura zgodna z~PEPPOL BIS Billing 3.0 jest z~definicji zgodna z~EN~16931, lecz nie odwrotnie. BIS 3.0 dodaje ponad 80~dodatkowych reguł walidacyjnych (artefakty ISO Schematron) zaostrza opcjonalność niektórych pól EN~16931 i wymaga konkretnych list kodów. Aktualne wydanie pochodzi z~listopada 2025~r. Specyfikacja obsługuje dwie składnie XML: UBL~2.1 (podstawowa) i CII D16B (alternatywna). \cite{peppol-bis-billing-3}

\paragraph{Porównanie z FA(3).}
Zestawiając trzy formaty, FA(3) i FatturaPA reprezentują podejście \textit{krajowe} -- oba są autonomicznymi schematami XML niepowiązanymi bezpośrednio z~EN~16931, definiowanymi przez administracje skarbowe wyłącznie na potrzeby własnego systemu. PEPPOL BIS~3.0 reprezentuje podejście \textit{europejskie} -- stanowi formalną implementację normy transgranicznej, co umożliwia wymianę faktur bez konwersji formatu między 40+~jurysdykcjami. Szczegółową strukturę FA(3) opisano w~sekcji~\ref{sec:ksef-wzor-faktury}.

\subsection{Modele komunikacji}
\label{sec:porownanie-modele-komunikacji}

\subsubsection{Protokoły transportowe i mechanizmy komunikacji}
\label{sec:porownanie-protokoly-transportowe-mechanizmy-komunikacji}

Trzy analizowane systemy różnią się zasadniczo w warstwie transportowej -- zarówno pod względem liczby obsługiwanych kanałów, jak i paradygmatu wymiany komunikatów.

\textbf{SdI} udostępnia cztery równoległe kanały transmisji: \cite{sdi-fatturapa-messages}
\begin{itemize}
    \item \textbf{SDICoop} -- kanał SOAP/Web Services przez Internet; wymaga akredytacji i~certyfikatu kanałowego.
    \item \textbf{SDIFTP} -- transfer plików FTP/SFTP na dedykowanych terminalach; wymaga akredytacji i stałej infrastruktury.
    \item \textbf{PEC} (\textit{Posta Elettronica Certificata}) -- włoska certyfikowana poczta e-mail; faktury wysyłane jako załącznik XML na wskazany adres SdI; nie wymaga akredytacji.
    \item \textbf{Interfejs webowy} ``Fatture e Corrispettivi'' -- upload pliku przez przeglądarkę (limit 5~MB); przeznaczony dla sporadycznych wystawców.
\end{itemize}

\textbf{PEPPOL} stosuje pojedynczy, ustandaryzowany protokół transportowy: \textbf{AS4} w profilu Peppol AS4 v2.0.3, opartym na specyfikacji OASIS ebMS~3.0 i~zgodnym z~eDelivery AS4 v1.14 Komisji Europejskiej. Wzorzec wymiany to \textit{one-way push} -- C2 wysyła wiadomość do C3 bez oczekiwania na odpowiedź biznesową w~tej samej sesji. Komunikat SOAP~1.2 zawiera podpis XML-DSig i szyfrowanie XML-Enc \textit{payloadu}; warstwa transportu wymaga TLS~1.2+. \cite{peppol-as4-profile,ec-as4-edelivery}

Kluczowym elementem infrastruktury PEPPOL jest mechanizm \textbf{SML/SMP}: nadawca oblicza hash identyfikatora PEPPOL odbiorcy i odpytuje \textbf{SML} (\textit{Service Metadata Locator}) -- serwer DNS wskazujący właściwy \textbf{SMP} (\textit{Service Metadata Publisher}). SMP zwraca URL endpointu AS4 C3 oraz jego klucz publiczny, umożliwiając dynamiczne wyszukiwanie bez centralnego rejestru adresów. \cite{peppol-sml-spec}

\textbf{KSeF} stosuje inny paradygmat: jednolity \textbf{REST API} opisany specyfikacją OpenAPI~3.0.4 (\ref{sec:ksef-interfejs-programistyczny}). W~odróżnieniu od asynchronicznego modelu SOAP/AS4, REST API KSeF jest synchroniczne w warstwie pojedynczego żądania -- klient wysyła zaszyfrowaną fakturę i~natychmiast otrzymuje identyfikator sesji. Brak mechanizmu dynamicznego wyszukiwania: routing wynika z~centralnej architektury huba (NIP odbiorcy w~treści faktury wystarczy do jednoznacznego adresowania).

\subsubsection{Uwierzytelnianie i autoryzacja}
\label{sec:porownanie-uwierzytelnianie-autoryzacja}

Trzy systemy stosują odmienne modele tożsamości, odzwierciedlające ich architektoniczne paradygmaty.

\textbf{SdI} wymaga \textbf{kwalifikowanego podpisu elektronicznego} na samej fakturze: akceptowane są formaty \textbf{XAdES-BES} (ETSI TS~101~903 v1.4.1, metoda \textit{enveloped}) lub \textbf{CAdES-BES} (ETSI TS~101~733 V1.7.4). Podpis musi być złożony certyfikatem kwalifikowanym wydanym po weryfikacji tożsamości wystawcy. Dla kanałów akredytowanych (SDICoop, SDIFTP) wymagany jest dodatkowo certyfikat kanałowy uwierzytelniający połączenie techniczne. \cite{sdi-fatturapa-messages}

\textbf{PEPPOL} nie wymaga podpisu kwalifikowanego \emph{na treści faktury biznesowej} -- integralność i~autentyczność zapewnia warstwa transportowa AS4 (XML-DSig). Access Pointy operują w~ramach hierarchicznej \textbf{PKI OpenPeppol} generacji~G3, obejmującej osobne Root~CA, Access~Point~CA i~SMP~CA dla środowisk testowego i~produkcyjnego (certyfikaty G3 obowiązkowe od 1~stycznia 2026~r.). Akredytacja nowego Access Pointa wymaga: (1)~członkostwa w~OpenPeppol lub zawarcia umowy Service Provider Agreement z~krajową Peppol Authority, (2)~zaliczenia testów na OpenPeppol Testbed, (3)~uzyskania certyfikatu produkcyjnego. \cite{peppol-pki-2025}

\textbf{KSeF} stosuje model opisany w~sekcji~\ref{sec:ksef-uwierzytelnianie-i-autoryzacja}: uwierzytelnianie oparte na tokenach JWT z procesem Challenge-Response, w~którym podpis XAdES lub zaszyfrowany token systemowy KSeF służy do jednorazowego uzyskania pary \texttt{accessToken}/\texttt{refreshToken}. Brak wymogu kwalifikowanego podpisu na treści dokumentu XML faktury -- tę rolę pełni podpis przy uwierzytelnianiu i obowiązkowe szyfrowanie \textit{envelope}.

\subsubsection{Bezpieczeństwo i szyfrowanie}
\label{sec:porownanie-bezpieczenstwo-szyfrowanie}

Trzy systemy realizują odmienne strategie kryptograficznej ochrony faktury:

\begin{itemize}
    \item \textbf{SdI} opiera bezpieczeństwo na \textbf{podpisie kwalifikowanym na dokumencie} (XAdES/CAdES). Treść faktury XML nie jest szyfrowana -- SdI nie wymaga \textit{envelope encryption} zawartości biznesowej. Transport zabezpieczony jest przez TLS/HTTPS lub dedykowany szyfrowany kanał SFTP. Integralność i niezaprzeczalność gwarantuje podpis kwalifikowany, który w~prawie unijnym ma skutek prawny równoważny podpisowi odręcznemu. \cite{sdi-fatturapa-messages}

    \item \textbf{PEPPOL} realizuje bezpieczeństwo na \textbf{poziomie wiadomości AS4}: profil Peppol AS4 wymaga podpisu XML-DSig (\textit{Binary Security Token}) oraz szyfrowania XML-Enc \textit{payloadu} -- tzw.~\textit{message-level encryption} między C2 a~C3. Transport wymaga TLS~1.2+. Na poziomie faktury biznesowej brak wymogu kwalifikowanego podpisu (krajowe Peppol Authority mogą go wprowadzić krajowo). \cite{peppol-as4-profile}

    \item \textbf{KSeF} łączy obie warstwy: TLS w~transporcie oraz obowiązkowe \textit{envelope encryption} każdej faktury (AES-256-CBC + RSA-OAEP) \emph{przed} przesłaniem do API -- zarówno w~sesji interaktywnej, jak i~wsadowej (\ref{sec:ksef-szyfrowanie-i-podpis}). Brak podpisu kwalifikowanego na dokumencie XML; integralność zapewnia: szyfrowanie \textit{envelope}, weryfikacja skrótów SHA-256 w~sesji i~podpis XAdES przy uwierzytelnianiu klienta. \cite{ksef-docs-przeglad,ksef-weryfikacja-faktury}
\end{itemize}

Podsumowując różnice strategii: SdI ufa podpisowi \emph{na dokumencie}, PEPPOL ufa zabezpieczeniom \emph{warstwy transportowej} AS4, a~KSeF stosuje obligatoryjne szyfrowanie \emph{zawartości} niezależnie od zabezpieczeń transportu.

\subsection{Podsumowanie}
\label{sec:porownanie-podsumowanie}

Analiza trzech systemów uwidacznia dwie fundamentalne linie podziału w~projektowaniu infrastruktury e-fakturowania. Pierwsza dotyczy \textbf{modelu architektonicznego}: SdI i~KSeF przyjęły podejście centralnego huba (\textit{clearance model}), w~którym państwowy pośrednik gwarantuje kompletność danych podatkowych w~czasie rzeczywistym. PEPPOL realizuje model sieci interoperacyjnej, w~którym zaufanie wynika z~akredytacji uczestników, a~nie z~centralnej kontroli dokumentów.

Druga linia podziału dotyczy \textbf{podejścia do standaryzacji formatu}. FatturaPA i~FA(3) to autonomiczne schematy XML definiowane przez krajowe administracje, niekompatybilne natywnie z~europejską normą EN~16931. PEPPOL BIS Billing~3.0 jest formalną implementacją tej normy, co umożliwia transgraniczną interoperacyjność bez konwersji formatu. Rodzi to pytanie o~kierunek harmonizacji europejskiej: czy krajowe systemy clearance'owe jak KSeF i~SdI powinny stopniowo adoptować EN~16931, czy też PEPPOL jest właściwą warstwą transgraniczną niezależną od krajowych platform?

\begin{xltabular}{\linewidth}{l X X X}
    \toprule
    \textbf{Aspekt} & \textbf{SdI (Włochy)} & \textbf{PEPPOL} & \textbf{KSeF (Polska)} \\
    \midrule
    \endfirsthead
    \toprule
    \textbf{Aspekt} & \textbf{SdI (Włochy)} & \textbf{PEPPOL} & \textbf{KSeF (Polska)} \\
    \midrule
    \endhead
    \bottomrule
    \endfoot
    \bottomrule
    \caption{Porównanie systemów e-fakturowania: SdI, PEPPOL i KSeF.}
    \label{tab:porownanie}
    \endlastfoot
    Operator & Agenzia delle Entrate / Sogei & OpenPeppol AISBL + krajowe Peppol Authorities & Ministerstwo Finansów RP \\
    Model architektoniczny & Central hub (clearance) & 4-corner / 5-corner (sieć interoperacyjna) & Central hub (clearance) \\
    Archiwum centralne & Nie & Nie & Tak (10 lat) \\
    Zakres podmiotowy & B2G + B2B + B2C & B2G (obowiązkowo UE), B2B (krajowo) & B2B obowiązkowo \\
    Format danych & FatturaPA XML (krajowy) & UBL 2.1 / CII D16B (PEPPOL BIS 3.0 / EN 16931) & FA(3) XML (krajowy) \\
    Zgodność z EN 16931 & Równoważna, nie natywna & Natywna (CIUS) & Nie \\
    Protokół transportowy & SDICoop (SOAP), SDIFTP, PEC, web UI & AS4 (ebMS 3.0), one-way push & REST API (HTTPS + JSON/XML) \\
    Discovery odbiorcy & Codice Destinatario / PEC w fakturze & SML (DNS) + SMP (dynamiczny lookup) & Brak (NIP w treści, centralny hub) \\
    Identyfikator podmiotu & Partita IVA + Codice Destinatario (7 znaków) & Peppol Participant ID (\texttt{ICD:value}) & NIP \\
    Podpis na fakturze & Obowiązkowy XAdES-BES lub CAdES-BES & Brak (opcjonalny krajowo) & Brak \\
    Szyfrowanie zawartości & Nie wymagane & XML-Enc \textit{payloadu} AS4 & AES-256-CBC + RSA-OAEP (obowiązkowe) \\
    Komunikat zwrotny & 7 typów \textit{notifiche} SdI & MLR (Accepted / Rejected) & UPO (XML) \\
    Specyfikacja publiczna & XSD + dok.\ techniczna fatturapa.gov.it & OpenAPI + Schematron na docs.peppol.eu & OpenAPI 3.0.4 + XSD FA(3) \\
\end{xltabular}
